\labsection{4}{Explore distributions, subgroups, and relationships}{sec:lab04}

\begin{labproject}{Exploratory analysis: develop a defensible EDA narrative}
\textbf{Coverage.} Chapter~\ref{ch:eda}.\\
\textbf{Goal.} You should be able to combine robust summaries and graphics, investigate subgroup structure, detect when aggregate relationships are misleading, and translate exploration into testable modelling questions.\\
\textbf{Dataset.} \texttt{medical\_cost\_personal\_dataset.csv}.\\
\textbf{Scenario.} An insurer wants to understand the structure of medical charges before any predictive model is considered. Your task is to discover patterns, identify plausible interactions, and separate observations that are unusual from observations that are invalid.

\courselabflow{Inspect charges}{Compare groups}{Check relationships}{Explain limits}


\textbf{Run the complete script.} Start a clean R session in the course folder. Each step continues the same file; comment headings identify the operation performed by the next block.
\begin{lstlisting}[style=utpterminal]
Rscript all_experiments/lab04_eda.R
\end{lstlisting}

\noindent\textbf{Step 1.} Load, validate and prepare the data
\lstinputlisting[style=utpr,firstline=1,lastline=8]{all_experiments/lab04_eda.R}

\noindent\textbf{Step 2.} 1. Target shape and robust summaries.
\courselabtrace{An outlier is a question, not a deletion rule}{Flag}{Use a stated threshold to identify unusual charges}{Inspect}{Compare age, BMI and smoking group}{Retain}{Remove only with a documented data-quality reason}



\textbf{Use the chapter while coding.} Revisit \hyperref[ch:eda]{the earlier chapter explanation}; match its operation to the variables and state changes below.

\textbf{Follow the supplied implementation.} The threshold is a descriptive rule, not proof of bad data. Compare the smoker and non-smoker summaries and inspect the original rows before proposing any exclusion. Report how many observations the flag captures.
\courselabtrace{Trace the supplied outlier flag}{1. Input / state}{q stores charge quartiles; iqr is IQR(d\$charges).}{2. Operation}{upper\_fence = q[3] + 1.5 times iqr; flagged keeps rows above that boundary.}{3. Check}{Inspect flagged records; the script does not remove them before subgroup comparisons.}


\lstinputlisting[style=utpr,firstline=9,lastline=19]{all_experiments/lab04_eda.R}

\noindent\textbf{Step 3.} Inspect flagged rows rather than deleting them.
\lstinputlisting[style=utpr,firstline=20,lastline=26]{all_experiments/lab04_eda.R}

\noindent\textbf{Step 4.} 2. Subgroup comparison.
\lstinputlisting[style=utpr,firstline=27,lastline=33]{all_experiments/lab04_eda.R}

\noindent\textbf{Step 5.} 3. Region is inspected without over-interpreting small differences.
\lstinputlisting[style=utpr,firstline=34,lastline=41]{all_experiments/lab04_eda.R}

\noindent\textbf{Step 6.} 4. Overall versus conditional relationship.
\lstinputlisting[style=utpr,firstline=42,lastline=60]{all_experiments/lab04_eda.R}

\noindent\textbf{Step 7.} Extension: age bands within smoking groups.
\lstinputlisting[style=utpr,firstline=61,lastline=65]{all_experiments/lab04_eda.R}
\textbf{Chapter connection.} Revisit \hyperref[ch:eda]{the chapter explanation} and map its states to the steps below.
% Reused from chapters/04-exploratory-analysis.tex
\begin{labvisual}{Recall the chapter model}
\begin{labgraphic}
\begin{tikzpicture}[x=\linewidth,y=1cm,
  lens/.style={draw=UTPBlue!65,fill=UTPPaleBlue,rounded corners=1.8mm,
    line width=0.7pt,text width=0.235\linewidth,
    minimum width=0.29\linewidth,minimum height=2.20cm,
    inner xsep=4mm,inner ysep=3mm,align=center,
    font=\scriptsize\color{UTPNavy}},
  flow/.style={-{Stealth[length=2mm,width=1.5mm]},
    line width=0.9pt,draw=UTPTeal}]
\node[lens] (u) at (0.17,0)
  {\textcolor{UTPBlue}{\textbf{01  DESCRIBE}}\\[5pt]
   {\footnotesize\textbf{UNIVARIATE}}\\[4pt]
   \textbf{One variable}\\[5pt]
   Shape, centre, and spread};
\node[lens] (b) at (0.50,0)
  {\textcolor{UTPBlue}{\textbf{02  RELATE}}\\[5pt]
   {\footnotesize\textbf{BIVARIATE}}\\[4pt]
   \textbf{Two variables}\\[5pt]
   Differences and relationships};
\node[lens,draw=UTPTeal!80,fill=UTPPaleTeal] (m) at (0.83,0)
  {\textcolor{UTPTeal}{\textbf{03  CONDITION}}\\[5pt]
   {\footnotesize\textbf{MULTIVARIATE}}\\[4pt]
   \textbf{Multiple variables}\\[5pt]
   Patterns within groups\\and joint structure};
\draw[flow] (u.east) -- (b.west);
\draw[flow] (b.east) -- (m.west);
\end{tikzpicture}
\end{labgraphic}\end{labvisual}


\begin{verification}
\textbf{Explain and extend your result.} Inspect one unusually high charge against similar-age records in the same smoking group. Explain what is unusual and what missing information prevents a causal explanation.


\textbf{Required evidence.} Submit an EDA brief containing three graphics and no more than six evidence statements. Each statement must include the comparison or relationship, its direction/magnitude, and an appropriate limitation.

\textbf{Professional standard.} Exploration is used to generate and refine hypotheses. Do not present a pattern discovered in the same data as if it were independently confirmed evidence.
\end{verification}
\end{labproject}
